% chktex-file 0
% ===== 序言 =====
\chapter{序言}

\textit{一个国家若不了解自己的过去，便如同一个失去记忆的人——他可以继续行走、继续呼吸，却永远无法回答那个最简单也最沉重的问题：我是谁？这片我们称之为故土的土地，承载了太多的荣耀与创伤，太多的欢歌与悲鸣。然而令人痛心的是，我们之中有太多的人，对自己脚下这片土地的过往，所知甚少。}

\textit{正是基于这样一种忧虑，也基于一种无法推卸的责任感，我写下了这本书。本书是出于教育与科普目的而作，旨在为那些对这块苦难的土地尚不了解的人，提供一个参考和快速入门。当然，也便于外邦的友人，对本国的历史情况，有一个大致的了解。如果能借此培养他们对本国历史的兴趣，便再好不过了——因为唯有了解，才能产生理解；唯有理解，才能孕育尊重。}

\textit{本书的内容，主要是根据现有的史料，结合作者的研究与筛选整理而成。这本小册子对这块土地上错综复杂的历史情况，提供一个大致的概览。然而我也必须坦率地告诉读者：这远非一部完整的叙事。经过谨慎的思考和选择之后，我们并不会为所有的篇章都提供详细的叙述——那些渗入这片土地的血与泪水，便留给读者自己去体会。有些伤痛，语言无力承载；有些沉默，比千言万语更有分量。}

\textit{这种取舍并非逃避。恰恰相反，我们不应该在每一个伤口上反复摩挲，我们应该让后人知晓这些伤口的存在，同时又赋予他们向前看的力量。铭记，但不沉溺；审视，但不绝望。我希望这本小书，能够成为这样一座桥梁——连接过去与现在，连接苦难与希望，连接每一个Eullossetch人与这片土地之间，那条看不见却从未断裂的纽带。}

\textit{显然，由于作者的水平有限，谬误和纰漏在所难免。我真诚地希望读者能及时反馈，为我们提供宝贵的建议。一部史学著作，从来不是一个人关在书斋里完成的——它需要在无数双眼睛的审视下，在无数种声音的碰撞中，逐渐趋向于真实。}

\textit{最后，感谢所有为本书提供帮助的人们。感谢 A.~Mikov，他为本文提供了宝贵的建议；感谢编辑 K.~Hydalavic，没有他的校对，本书是不可能完成的。以及所有在这个过程中帮助过我的人，包括那些志愿帮我收集资料的学生们。}

\textit{谨以此书，献给所有在这片土地上生活过、奋斗过、痛苦过、也欢笑过的人们。}

\textit{写于\blackout{1978/08/15}，南丘卡拉大学。}
